\documentclass[../report.tex]{subfiles}
\begin{document}

\section{Trạng thái kiến trúc và định dạng lệnh}\label{sec:isa}

RV32I có 32 thanh ghi x0 đến x31, mỗi thanh ghi 32 bit, trong đó x0 luôn đọc ra 0 và mọi lệnh ghi vào x0 bị bỏ
qua \cite{riscv-spec}. Vì x0 là hằng 0, phần cứng chỉ cần 31 thanh ghi thật x1 đến x31
(mục~\ref{sec:regfile}). Bộ nhớ đánh địa chỉ theo byte, kiểu little-endian. Mỗi lệnh dài 4 byte nên PC tăng 4 sau
mỗi lệnh. ROM và RAM của Logisim lưu mỗi ô một word, nên chỉ số ô bằng địa chỉ byte bỏ 2 bit thấp: \sig{IMEM}
dùng \sig{PC[11:2]}, \sig{DMEM} dùng \sig{Addr[13:2]}.

Mọi lệnh có opcode ở bit [6:0]. Các trường rd [11:7], funct3 [14:12], rs1 [19:15], rs2 [24:20] luôn ở cùng vị trí
trong mọi định dạng có trường đó (bảng~\ref{tab:fmt}). Nhờ vậy register file đọc rs1, rs2 song song với việc giải
mã. Bit dấu của immediate luôn là bit 31, nên phép mở rộng dấu bắt đầu được ngay mà không cần biết định dạng.

\begin{table}[H]
\centering\small
\caption{Sáu định dạng lệnh RV32I}\label{tab:fmt}
\begin{tabular}{@{}lllllll@{}}
\toprule
Định dạng & [31:25] & [24:20] & [19:15] & [14:12] & [11:7] & [6:0] \\
\midrule
R & funct7 & rs2 & rs1 & funct3 & rd & opcode \\
I & \multicolumn{2}{l}{imm[11:0]} & rs1 & funct3 & rd & opcode \\
S & imm[11:5] & rs2 & rs1 & funct3 & imm[4:0] & opcode \\
B & imm[12|10:5] & rs2 & rs1 & funct3 & imm[4:1|11] & opcode \\
U & \multicolumn{4}{l}{imm[31:12]} & rd & opcode \\
J & \multicolumn{4}{l}{imm[20|10:1|11|19:12]} & rd & opcode \\
\bottomrule
\end{tabular}
\end{table}

Hai CPU hiện thực 37 lệnh RV32I: 10 lệnh R, 9 lệnh I số học, 5 lệnh load, 3 lệnh store, 6 lệnh rẽ nhánh có
điều kiện, JAL, JALR, LUI và AUIPC. FENCE được giải mã thành NOP: \sig{Control} không có nhóm cho opcode
\texttt{0001111}, nên mọi tín hiệu điều khiển bằng 0. Cách này đúng với hai CPU ở đây vì chúng chỉ có một luồng
thực thi, không có cache, và mọi truy cập RAM hay I/O diễn ra lần lượt theo đúng thứ tự chương trình, nên FENCE
không còn thứ tự nào cần ép thêm. Các lệnh CSR thuộc phần mở rộng Zicsr, không nằm trong RV32I base, và không được
hiện thực. ECALL và EBREAK được quy ước là lệnh dừng CPU, để
bộ test có điểm kết thúc rõ ràng.

\section{Sơ đồ khối và các quyết định thiết kế}

Datapath gồm thanh ghi \sig{PC}, ROM \sig{IMEM}, khối giải mã \sig{Control} và \sig{ImmGen}, \sig{RegisterFile},
\sig{ALU} cùng \sig{BranchUnit}, RAM \sig{DMEM} với \sig{LoadUnit}, \sig{StoreUnit}, và MUX chọn giá trị ghi lại.
Các circuit con này dùng chung cho SC và PL; PL thêm 4 thanh ghi pipeline và 3 khối xử lý hazard. Bảng~\ref{tab:dec}
ghi các quyết định thiết kế chính.

\begin{table}[H]
\centering\small
\caption{Các quyết định thiết kế}\label{tab:dec}
\begin{tabularx}{\textwidth}{@{}>{\raggedright\arraybackslash}p{4.3cm}X@{}}
\toprule
Quyết định & Lý do và đánh đổi \\
\midrule
Kiến trúc Harvard: \sig{IMEM} là ROM, \sig{DMEM} là RAM riêng & Nạp lệnh và load/store không tranh nhau một cổng bộ nhớ, tránh structural hazard giữa tầng IF và MEM. Đổi lại chương trình không tự sửa mã được. \\
ALU tính luôn địa chỉ đích & Với branch và JAL, đầu vào A của ALU là PC; với JALR là rs1. Đích nhảy = kết quả ALU AND \texttt{0xFFFFFFFE}. Không cần bộ cộng PC + imm riêng; phép so sánh rs1 với rs2 vì thế làm ở \sig{BranchUnit}. \\
Lệnh \texttt{0x00000000} là NOP & Decoder trong \sig{Control} chỉ bật khi \sig{opcode[1:0]} = 11, nên lệnh toàn 0 cho mọi tín hiệu điều khiển bằng 0. Flush một tầng pipeline chỉ cần nạp 0. \\
SB, SH bằng đọc, sửa rồi ghi cả word & RAM của Logisim chỉ có byte enable khi đọc đồng bộ, không dùng được cho CPU đơn chu kỳ. \\
Memory-mapped I/O & UART và GPIO là vài địa chỉ đặc biệt; CPU dùng chính \texttt{lw}, \texttt{sw}, \texttt{lb}, \texttt{sb} để điều khiển, không cần lệnh mới. \\
\bottomrule
\end{tabularx}
\end{table}

\section{Bản đồ bộ nhớ}

\begin{table}[H]
\centering\small
\caption{Bản đồ bộ nhớ}\label{tab:memmap}
\begin{tabularx}{\textwidth}{@{}ll>{\raggedright\arraybackslash}X@{}}
\toprule
Vùng & Địa chỉ & Ghi chú \\
\midrule
\sig{IMEM} (ROM) & \texttt{0x0000\_0000}--\texttt{0x0000\_0FFF} & 1024 lệnh, địa chỉ ô = \sig{PC[11:2]} \\
\sig{DMEM} (RAM) & \texttt{0x0000\_0000}--\texttt{0x0000\_3FFF} & 16 KiB, địa chỉ ô = \sig{Addr[13:2]} \\
Vùng dump của bộ test & \texttt{0x0000\_0700}--\texttt{0x0000\_077F} & \file{run\_tests.py} lưu x1 đến x31 vào \texttt{0x704}--\texttt{0x77C} trước khi dừng; ô \texttt{0x700} ứng với x0 để trống \\
UART\_TXDATA & \texttt{0xF000\_0000} (ghi) & ký tự ra TTY \\
UART\_RXDATA & \texttt{0xF000\_0004} (đọc) & lấy một ký tự từ Keyboard \\
UART\_STATUS & \texttt{0xF000\_0008} (đọc) & bit 0: có ký tự chờ đọc; bit 1: TX sẵn sàng \\
GPIO\_OUT & \texttt{0xF000\_0010} (đọc/ghi) & bit [15:0] ra 4 LED 7 đoạn, bit [23:16] ra 8 LED \\
GPIO\_IN & \texttt{0xF000\_0014} (đọc) & 8 công tắc DIP \\
\bottomrule
\end{tabularx}
\end{table}

Thanh ghi và PC reset về 0; phần mềm tự đặt \texttt{sp}, các chương trình test dùng \texttt{sp = 0x3FF0}.

\end{document}
